\section{Cuaderno de Google Colab}
    Todo el desarrollo descrito en este informe se implementó en un único
    cuaderno de Google Colab titulado \textit{Predicción del puntaje global de la
    prueba Saber 11 a partir de variables socioeconómicas y del colegio}, que se
    conserva también en el repositorio del proyecto como
    \texttt{notebooks/01\_planteamiento\_y\_eda.ipynb}. Los segmentos de código
    reproducidos en las secciones de metodología y resultados corresponden, uno a
    uno y en el mismo orden, a las celdas de ese cuaderno, y cada uno tiene su
    archivo espejo en \texttt{utils/codes/}.

    \vspace{10pt}

    \begin{tcolorbox}[
        sharp corners,
        colback=green!7!white,
        colframe=green!55!black,
        boxrule=1.2pt,
        fonttitle=\bfseries,
        title=Cuaderno interactivo disponible en Google Colab
    ]
        \centering

        \hypersetup{pdfborder={0 0 0}}
        \urlstyle{same}

        {\color{green!40!black}%
        \url{https://colab.research.google.com/drive/120GLTvDurcx6kycc91pffk8ilKFocwyq?usp=sharing}}

    \end{tcolorbox}

    \vspace{10pt}

    El cuaderno está organizado en doce secciones que siguen el mismo hilo
    argumental del informe:

    \vspace{10pt}

    \begin{itemize}
        \item Planteamiento del problema, pregunta de investigación, objetivos,
        hipótesis y alcance.

        \item Preparación del entorno e identidad gráfica.

        \item Carga y descripción del conjunto de datos, con el diccionario de
        variables.

        \item Estructura y calidad de los datos: duplicados, columnas constantes,
        valores faltantes y conversión de tipos.

        \item La variable objetivo: distribución y verificación de la fórmula del
        puntaje global.

        \item Variables del estudiante.

        \item Variables del hogar.

        \item Variables del colegio.

        \item Variables geográficas.

        \item Correlaciones y ajuste lineal exploratorio.

        \item Hallazgos y decisiones para la fase 2.

        \item Conclusiones.
    \end{itemize}

    \vspace{10pt}

    Su ejecución no requiere ninguna instalación ni descarga manual, puesto que el
    conjunto de datos se obtiene con \textit{kagglehub} desde la copia pública
    publicada en Kaggle \cite{rubiano2026kaggle}, sin necesidad de cuenta, y todas
    las dependencias vienen preinstaladas en el entorno estándar de Colab. La
    descarga del archivo Parquet de 31,5 MB toma unos segundos y el cuaderno
    completo se ejecuta en menos de cinco minutos. La única fuente de aleatoriedad es el
    muestreo del ajuste lineal exploratorio, gobernado por la semilla
    \texttt{RANDOM\_STATE = 42}, de modo que todas las cifras reportadas se
    reproducen exactamente.
